\documentclass[11pt]{article}
\usepackage[T1]{fontenc}
\usepackage[utf8]{inputenc}
\usepackage{amsmath,amssymb,amsthm,mathtools}
\usepackage[margin=1in]{geometry}
\setlength{\emergencystretch}{1em}
\newtheorem{theorem}{Theorem}
\newtheorem{lemma}{Lemma}

\title{A Disproof of Conjecture 00000000154}
\author{Prepared for the account \texttt{idealistichacker}}
\date{}

\begin{document}
\maketitle

\begin{abstract}
The conjecture claims that infinitely many even-indexed derangement numbers
are prime. We prove the stronger result that no even-indexed derangement
number is prime. A separate Lean~4 project formalizes the statement on actual
fixed-point-free permutations and proves that the relevant index set is empty.
\end{abstract}

\section*{Frozen statement and scope}
The organizer source is identified below and was frozen before this analysis:
\begin{flushleft}
\footnotesize
Source: \texttt{conjectures/00000000154.md}\\
Commit: \texttt{95acb520ec5607c826b8a997b1ef2fc82d6f7c57}
\end{flushleft}
It states:

\begin{quote}
``Prime values of derangement numbers: there exist infinitely many even $n$
such that $D_n$ is prime ($D_n$ can be odd only when $n$ is even).''
\end{quote}

\begin{flushleft}
\footnotesize
Git blob: \texttt{5d2fb2e814a4c521d2ace36034863cc8ec6fe836}\\
SHA-256: \texttt{694c389d34ad78b41caae8ad4e1f414fa0541e6b19b81aeab432466039dc87c1}
\end{flushleft}
Here $D_n$ has its standard meaning: the number of fixed-point-free
permutations of an $n$-element set.

The parenthetical observation about when $D_n$ can be odd is not a hypothesis
or a conclusion of the present proof.  We prove only the negation of the
source's asserted infinite set of \emph{even} indices with prime $D_n$.

\section*{Derangement recurrence}
The standard initial values and recurrence are
\[
 D_0=1,\qquad D_1=0,\qquad
 D_{r+2}=(r+1)\bigl(D_{r+1}+D_r\bigr) \quad (r\geq0).
\]
In particular, $D_2=(1)(D_1+D_0)=1$.

\begin{lemma}
For every $r\geq0$, $D_{r+2}>0$.
\end{lemma}
\begin{proof}
Induct on $r$.  The base case is $D_2=1>0$.  If $D_{r+2}>0$, then
\[
D_{r+3}=(r+2)\bigl(D_{r+2}+D_{r+1}\bigr)>0,
\]
since $r+2>0$, $D_{r+2}>0$, and derangement numbers are nonnegative.
\end{proof}

\begin{theorem}
For every even natural number $n$, $D_n$ is not prime.
\end{theorem}
\begin{proof}
Write $n=2k$.  There are three branches.

If $k=0$, then $n=0$ and $D_n=D_0=1$, not prime.  If $k=1$, then $n=2$ and
$D_n=D_2=1$, also not prime.  Finally suppose $k\geq2$, so $n\geq4$.
The recurrence gives
\[
 D_n=(n-1)\bigl(D_{n-1}+D_{n-2}\bigr).
\]
The first factor obeys $n-1\geq3$.  Also $n-2\geq2$ and $n-1\geq2$, so the
lemma gives $D_{n-2}>0$ and $D_{n-1}>0$.  Thus
$D_{n-1}+D_{n-2}\geq2$.  This expresses $D_n$ as a product of two natural
numbers both greater than one; it cannot be prime.
\end{proof}

Thus
\[
 \{n\in\mathbb N : n\text{ is even and }D_n\text{ is prime}\}=\varnothing,
\]
which is certainly not infinite.  This disproves the conjecture.

\section*{Lean statement bridge}
The Lean source does not formalize an unrelated recurrence alone.  Mathlib's
\texttt{derangements (Fin n)} is the subtype of fixed-point-free permutations
of \texttt{Fin n}; \texttt{Fintype.card} counts that actual subtype.  The
Mathlib theorem \texttt{card\_derangements\_fin\_eq\_numDerangements} connects
this cardinality to the recurrence above.

The final Lean statements are:
\[
\forall n\in\mathbb N,\quad \mathrm{Even}(n)\to
\neg\mathrm{Prime}\bigl(\#\,\texttt{derangements (Fin n)}\bigr),
\]
and equality of the set
\[
\{n:\mathrm{Even}(n)\land
\mathrm{Prime}(\#\,\texttt{derangements (Fin n)})\}
\]
with the empty set, followed by its non-infinitude.  Hence the formalization
contains an explicit bridge from the all-even non-primality theorem to the
exact set-existence wording of the source conjecture.

\texttt{Check.lean} prints the transitive axiom dependencies of every
advertised theorem.  The elementary recurrence lemmas print
{\footnotesize\texttt{[propext, Quot.sound]}}; the statements about actual derangement
cardinalities, the all-even bridge, emptiness, and non-infinitude print
{\footnotesize\texttt{[propext, Classical.choice, Quot.sound]}}.  These standard Lean/Mathlib
principles are disclosed rather than called zero axioms; no custom axiom,
\texttt{sorry}, \texttt{admit}, \texttt{native\_decide}, or unsafe evaluation
is used.

\section*{Status}
This package includes a PDF compiled from this source with Tectonic~0.15.0
with deterministic PDF timestamp metadata. An internal independent-AI review
record is included in \texttt{review.json}.

This is a review request, not a claim of organizer approval, merge, acceptance,
ranking, award, or human verification. Any future revision must repeat the
mathematical and statement-alignment review, source/duplicate checks, and
reproducible PDF verification.

\end{document}
